\begin{tikzpicture}[font=\small]
  % samples used by all three panels
  \def\samps{{0.95,1.32,1.18,0.62,0.28,0.55,1.05,1.25}}

  % ---------- (a) zero-order hold ----------
  \begin{scope}
    \node[note, anchor=south] at (3.2,2.2){(a) 零阶保持（ZOH）};
    \draw[ax] (-0.4,0) -- (7.0,0); \node[note, anchor=north west] at (6.8,-0.05){$t$};
    \draw[ax] (0,-0.2) -- (0,1.9);
    \foreach \i in {0,...,7}{
      \pgfmathsetmacro\yv{\samps[\i]}
      \draw[curve2, line width=1.3pt] (\i*0.8,\yv) -- (\i*0.8+0.8,\yv);
      \ifnum\i<7
        \pgfmathsetmacro\ynext{\samps[\i+1]}
        \draw[curve2, line width=1.3pt] (\i*0.8+0.8,\yv) -- (\i*0.8+0.8,\ynext);
      \fi
      \filldraw[ssblue] (\i*0.8,\yv) circle (2.2pt);
    }
    \node[note, anchor=north, align=center] at (3.2,-0.4)
      {$h(t)$ 是宽度 $T$ 的矩形 $\Rightarrow$ $H(j\omega)$ 是 sinc，\\
       通带内下垂、带外泄漏都很大};
  \end{scope}

  % ---------- (b) ideal band-limited interpolation ----------
  \begin{scope}[xshift=9.4cm]
    \node[note, anchor=south] at (3.2,2.2){(b) 理想带限内插（sinc 内插）};
    \draw[ax] (-0.4,0) -- (7.0,0); \node[note, anchor=north west] at (6.8,-0.05){$t$};
    \draw[ax] (0,-0.2) -- (0,1.9);
    % smooth curve through the samples
    \draw[curve3, line width=1.4pt]
      (0,0.95) .. controls (0.27,1.10) and (0.53,1.26) .. (0.8,1.32)
      .. controls (1.07,1.38) and (1.33,1.28) .. (1.6,1.18)
      .. controls (1.87,1.08) and (2.13,0.78) .. (2.4,0.62)
      .. controls (2.67,0.46) and (2.93,0.28) .. (3.2,0.28)
      .. controls (3.47,0.28) and (3.73,0.42) .. (4.0,0.55)
      .. controls (4.27,0.68) and (4.53,0.95) .. (4.8,1.05)
      .. controls (5.07,1.15) and (5.33,1.24) .. (5.6,1.25);
    \foreach \i in {0,...,7}{
      \pgfmathsetmacro\yv{\samps[\i]}
      \draw[helper] (\i*0.8,0) -- (\i*0.8,\yv);
      \filldraw[ssblue] (\i*0.8,\yv) circle (2.2pt);
    }
    \node[note, anchor=north, align=center] at (3.2,-0.4)
      {$h(t)=\dfrac{\sin(\pi t/T)}{\pi t/T}$：在采样点处 $=1$，\\
       在其余采样点处 $=0$ $\Rightarrow$ 精确穿过每个样本};
  \end{scope}

  % ---------- frequency view ----------
  \begin{scope}[yshift=-4.6cm, xshift=4.4cm]
    \draw[ax] (-6.2,0) -- (6.6,0); \node[note, anchor=north west] at (6.4,-0.05){$\omega$};
    \draw[ax] (0,-0.2) -- (0,1.8);
    \foreach \c in {-4.0,0,4.0}{
      \draw[sslight, line width=1.2pt] (\c-1.2,0) -- (\c,1.35) -- (\c+1.2,0);
    }
    % ideal LP brick wall
    \draw[curve3, line width=1.4pt] (-2.0,0) -- (-2.0,1.55) -- (2.0,1.55) -- (2.0,0);
    % ZOH sinc envelope
    \draw[curve2, line width=1.3pt, domain=-6.0:6.0, samples=300, smooth]
      plot (\x, {1.55*abs(sin(deg(0.7854*\x+0.00001))/(0.7854*\x+0.00001))});
    \node[ssteal, font=\footnotesize, anchor=south east] at (-2.1,1.5){理想低通};
    \node[ssred,  font=\footnotesize, anchor=west] at (2.15,0.95){ZOH 的 sinc 包络};
    \node[note, anchor=north] at (4.0,-0.1){$\omega_s$};

    \node[note, anchor=north, align=center] at (0,-0.75)
      {ZOH 既\textbf{压低了通带}（下垂），又\textbf{漏进了副本}。
       实用做法：ZOH 之后再补一级模拟低通，并在数字端预加重补偿下垂};
  \end{scope}
\end{tikzpicture}
